\documentclass[aspectratio=169]{beamer}
\usepackage{pgfpages}
\setbeameroption{show notes on second screen=right}
\usetheme{Madrid}
\title{Reasons and Persons}
\subtitle{A Sample Dual-Screen Deck}
\author{R. Ridenour}
\institute{RRShow Test Fixture}
\date{\today}

\begin{document}

\begin{frame}
  \titlepage
\end{frame}
\note{Open by thanking the department. Mention that the deck was compiled with
\texttt{show notes on second screen=right}, which is exactly the shape RRShow
is built to split.}

\begin{frame}{What the Audience Sees}
  \begin{itemize}
    \item The left half of every page is the slide.
    \item It goes to the projector, full bleed.
    \item Nothing on this half is private.
  \end{itemize}
\end{frame}
\note{This is the notes half. It should never reach the external display.
\begin{itemize}
  \item Check that the crop is exact --- no sliver of notes bleeding into the slide.
  \item Check that the aspect ratio stays 16:9 after cropping.
\end{itemize}}

\begin{frame}{Three Arguments}
  \begin{enumerate}
    \item The appeal to parsimony
    \item The appeal to explanatory scope
    \item The appeal to intuition
  \end{enumerate}
  \vfill
  \begin{block}{Key claim}
    Personal identity is not what matters in survival.
  \end{block}
\end{frame}
\note{Spend the most time on the third argument. Expect pushback on the intuition
pump; have the teletransporter case ready.

Timing: aim to be here at the twenty minute mark.}

\begin{frame}{A Frame With No Note}
  Some frames simply have no note attached. The notes half of the page is then
  blank, and RRShow should render an empty pane rather than failing.
\end{frame}

\begin{frame}{Conclusion}
  \begin{itemize}
    \item Reductionism is defensible.
    \item The practical upshot is considerable.
  \end{itemize}
\end{frame}
\note{Close on the practical upshot. Take questions.

Remember to press B to black the screen during discussion.}

\end{document}
